\chapter{Operators}

This chapter develops the theory of linear operators between vector
lattices and Banach lattices. We begin with positive operators —
those linear maps which preserve the positive cone — establishing
their equivalence with monotonicity, the absolute-value estimate
$\lvert T x \rvert \le T \lvert x \rvert$, the natural partial
order on the space of linear operators, the extension of additive
maps from the positive cone into any Archimedean vector lattice,
and the automatic continuity of positive operators on Banach
lattices.

\section{Positive operators}

This section introduces positive operators between vector lattices
and records their first properties: the equivalence of positivity
and monotonicity, the absolute-value estimate $\lvert T x \rvert
\le T \lvert x \rvert$, and the natural partial order on the space
of linear operators. We then prove the extension lemma — every
additive map on the positive cone extends uniquely to a positive
operator into an Archimedean vector lattice — and we close with the
automatic continuity of positive operators on Banach lattices.

\begin{definition}
\label{def:positive-operator}
\lean{Positive}
\leanok
\uses{def:vector-lattice}
Let $X$ and $Y$ be vector lattices. A linear map $f \colon X \to Y$
is called \emph{positive} when $f x \ge 0$ for every $x \in X$
with $x \ge 0$.
\end{definition}

\begin{theorem}
\label{thm:positive-monotone-iff}
\lean{Positive.monotone_iff}
\leanok
\uses{def:positive-operator, def:vector-lattice}
Let $X$ and $Y$ be vector lattices and let $f \colon X \to Y$ be a
linear map. Then $f$ is monotone if and only if $f$ is positive.
\end{theorem}

\begin{theorem}
\label{thm:positive-abs-le-map-abs}
\lean{Positive.abs_le_map_abs}
\leanok
\uses{def:positive-operator, def:vector-lattice,
  thm:positive-monotone-iff}
Let $X$ and $Y$ be vector lattices and let $f \colon X \to Y$ be a
positive linear map. Then $\lvert f x \rvert \le f \lvert x \rvert$
for every $x \in X$.
\end{theorem}

\begin{lemma}
\label{lem:positive-le-iff}
\lean{Positive.le_iff}
\leanok
\uses{def:positive-operator, def:vector-lattice}
Let $X$ and $Y$ be vector lattices. The space of linear maps
$X \to Y$ is partially ordered by declaring $T \le S$ when $S - T$
is positive; equivalently, $T \le S$ holds if and only if $T x \le
S x$ for every $x \in X$ with $x \ge 0$.
\end{lemma}

\begin{lemma}
\label{lem:positive-zero-le-iff}
\lean{Positive.zero_le_iff}
\leanok
\uses{def:positive-operator, lem:positive-le-iff}
Let $X$ and $Y$ be vector lattices and let $T \colon X \to Y$ be a
linear map. Then $0 \le T$ in the operator order if and only if
$T$ is positive.
\end{lemma}

\begin{theorem}
\label{thm:positive-extension}
\lean{Positive.extension}
\leanok
\uses{def:positive-operator, def:vector-lattice,
  def:is-vl-archimedean}
Let $X$ be a vector lattice, let $Y$ be an Archimedean vector
lattice, and let $\tau \colon X \to Y$ be a map satisfying $\tau x
\ge 0$ for every $x \ge 0$ in $X$ and $\tau (x + y) = \tau x +
\tau y$ for every $x, y \ge 0$ in $X$. Then there exists a linear
map $T \colon X \to Y$ given by $T x = \tau x^{+} - \tau x^{-}$.
\end{theorem}

\begin{lemma}
\label{lem:positive-extension-nonneg}
\lean{Positive.extension_nonneg}
\leanok
\uses{thm:positive-extension}
Let $X$ be a vector lattice, let $Y$ be an Archimedean vector
lattice, and let $\tau \colon X \to Y$ satisfy the hypotheses of
the extension theorem. Then $T x = \tau x$ for every $x \in X$
with $x \ge 0$, where $T$ denotes the linear extension of $\tau$.
\end{lemma}

\begin{theorem}
\label{thm:positive-extension-positive}
\lean{Positive.extension_positive}
\leanok
\uses{def:positive-operator, thm:positive-extension,
  lem:positive-extension-nonneg}
Let $X$ be a vector lattice, let $Y$ be an Archimedean vector
lattice, and let $\tau \colon X \to Y$ satisfy the hypotheses of
the extension theorem. Then the linear extension $T$ of $\tau$ is
a positive operator.
\end{theorem}

\begin{theorem}
\label{thm:positive-extension-unique}
\lean{Positive.extension_unique}
\leanok
\uses{def:positive-operator, thm:positive-extension,
  lem:positive-extension-nonneg}
Let $X$ be a vector lattice, let $Y$ be an Archimedean vector
lattice, and let $\tau \colon X \to Y$ satisfy the hypotheses of
the extension theorem. If $f \colon X \to Y$ is any positive
linear map with $f x = \tau x$ for every $x \ge 0$ in $X$, then
$f$ coincides with the linear extension $T$ of $\tau$.
\end{theorem}

\begin{theorem}
\label{thm:positive-continuous}
\lean{Positive.continuous}
\leanok
\uses{def:positive-operator, def:banach-lattice,
  def:normed-vector-lattice, thm:positive-abs-le-map-abs,
  thm:positive-monotone-iff}
Let $X$ be a Banach lattice and let $Y$ be a normed vector
lattice. Every positive linear operator $f \colon X \to Y$ is
continuous.
\end{theorem}

\section{Order bounded operators}

A linear map between vector lattices is \emph{order bounded} if it
sends order bounded sets to order bounded sets — equivalently, for
every positive $x$ there is a positive $y$ dominating $\lvert T z
\rvert$ for every $z$ with $\lvert z \rvert \le x$. Every positive
operator is order bounded; the order bounded operators form a
partially ordered real vector space; and every order bounded
operator from a Banach lattice into a normed vector lattice is
automatically continuous.

\begin{definition}
\label{def:isOrderBounded}
\lean{IsOrderBounded}
\leanok
\uses{def:vector-lattice}
Let $X$ and $Y$ be vector lattices. A linear map $f \colon X \to
Y$ is called \emph{order bounded} if, for every $x \in X$ with $0
\le x$, there exists $y \in Y$ with $0 \le y$ such that $\lvert f
z \rvert \le y$ for every $z \in X$ with $\lvert z \rvert \le x$.
\end{definition}

\begin{theorem}
\label{thm:positive-isOrderBounded}
\lean{Positive.isOrderBounded}
\leanok
\uses{def:isOrderBounded, def:positive-operator,
  thm:positive-abs-le-map-abs, thm:positive-monotone-iff}
Let $X$ and $Y$ be vector lattices. Every positive linear operator
$f \colon X \to Y$ is order bounded.
\end{theorem}

\begin{lemma}
\label{lem:isOrderBounded-add}
\lean{IsOrderBounded.add}
\leanok
\uses{def:isOrderBounded}
Let $X$ and $Y$ be vector lattices and let $f, g \colon X \to Y$
be order bounded linear maps. Then $f + g$ is order bounded.
\end{lemma}

\begin{lemma}
\label{lem:isOrderBounded-neg}
\lean{IsOrderBounded.neg}
\leanok
\uses{def:isOrderBounded}
Let $X$ and $Y$ be vector lattices and let $f \colon X \to Y$ be
an order bounded linear map. Then $-f$ is order bounded.
\end{lemma}

\begin{lemma}
\label{lem:isOrderBounded-sub}
\lean{IsOrderBounded.sub}
\leanok
\uses{def:isOrderBounded, lem:isOrderBounded-add,
  lem:isOrderBounded-neg}
Let $X$ and $Y$ be vector lattices and let $f, g \colon X \to Y$
be order bounded linear maps. Then $f - g$ is order bounded.
\end{lemma}

\begin{lemma}
\label{lem:isOrderBounded-smul}
\lean{IsOrderBounded.smul}
\leanok
\uses{def:isOrderBounded}
Let $X$ and $Y$ be vector lattices, let $f \colon X \to Y$ be an
order bounded linear map, and let $r \in \mathbb{R}$. Then $r f$
is order bounded.
\end{lemma}

\begin{theorem}
\label{thm:isOrderBounded-continuous}
\lean{IsOrderBounded.continuous}
\leanok
\uses{def:isOrderBounded, def:banach-lattice,
  def:normed-vector-lattice}
Let $X$ be a Banach lattice and let $Y$ be a normed vector
lattice. Every order bounded linear operator $f \colon X \to Y$
is continuous.
\end{theorem}

\begin{definition}
\label{def:orderBoundedHom}
\lean{OrderBoundedHom}
\leanok
\uses{def:isOrderBounded}
Let $X$ and $Y$ be vector lattices. The space $\mathcal{L}_b(X,
Y)$ of \emph{order bounded operators} from $X$ to $Y$ is the set
of order bounded real-linear maps $X \to Y$. It carries the
pointwise real vector space structure, and the partial order
$f \le g$ defined by $f x \le g x$ for every $x \in X$ with $0
\le x$, under which $\mathcal{L}_b(X, Y)$ is a partially ordered
real vector space.
\end{definition}

\section{Regular operators}

A linear map between vector lattices is \emph{regular} if it can
be written as the difference of two positive operators. Every
regular operator is order bounded.

\begin{definition}
\label{def:isRegularOp}
\lean{IsRegularOp}
\leanok
\uses{def:positive-operator, def:vector-lattice}
Let $X$ and $Y$ be vector lattices. A linear map $f \colon X \to
Y$ is called \emph{regular} if there exist positive linear maps
$g, h \colon X \to Y$ such that $f = g - h$.
\end{definition}

\begin{theorem}
\label{thm:isRegularOp-isOrderBounded}
\lean{IsRegularOp.isOrderBounded}
\leanok
\uses{def:isRegularOp, def:isOrderBounded,
  thm:positive-isOrderBounded, lem:isOrderBounded-sub}
Let $X$ and $Y$ be vector lattices. Every regular linear operator
$f \colon X \to Y$ is order bounded.
\end{theorem}

\section{Vector lattice homomorphisms and isomorphisms}

A \emph{vector lattice homomorphism} between two vector lattices is
a real-linear map that preserves the lattice operations $\vee$ and
$\wedge$. Such maps are automatically monotone, preserve absolute
values, positive and negative parts, and disjointness; they admit
two convenient characterisations among linear maps — preservation
of the absolute value, and (in the positive case) preservation of
disjointness. We then introduce \emph{vector lattice isomorphisms}
— bijective vector lattice homomorphisms with vector lattice
inverse — and show that any additive bijection between positive
cones extends to a vector lattice isomorphism into an Archimedean
codomain. In the Banach lattice setting, every vector lattice
isomorphism is automatically a continuous linear equivalence. We
close with the more rigid notion of \emph{Banach lattice
isometry}: a real linear isometric equivalence which preserves
$\vee$ and $\wedge$.

\subsection{Vector lattice homomorphisms}

\begin{definition}
\label{def:vecLatHom}
\lean{VecLatHom}
\leanok
\uses{def:vector-lattice}
Let $X$ and $Y$ be vector lattices. A \emph{vector lattice
homomorphism} from $X$ to $Y$ is a real-linear map $f \colon X
\to Y$ such that $f(x \vee y) = f x \vee f y$ and $f(x \wedge y) =
f x \wedge f y$ for all $x, y \in X$.
\end{definition}

\begin{lemma}
\label{lem:vecLatHom-map-abs}
\lean{VecLatHom.map_abs}
\leanok
\uses{def:vecLatHom}
Let $X$ and $Y$ be vector lattices and let $f \colon X \to Y$ be a
vector lattice homomorphism. Then $f \lvert x \rvert = \lvert f x
\rvert$ for every $x \in X$.
\end{lemma}

\begin{lemma}
\label{lem:vecLatHom-monotone}
\lean{VecLatHom.monotone}
\leanok
\uses{def:vecLatHom}
Let $X$ and $Y$ be vector lattices. Every vector lattice
homomorphism $f \colon X \to Y$ is monotone.
\end{lemma}

\begin{lemma}
\label{lem:vecLatHom-map-nonneg}
\lean{VecLatHom.map_nonneg}
\leanok
\uses{def:vecLatHom, lem:vecLatHom-monotone}
Let $X$ and $Y$ be vector lattices, let $f \colon X \to Y$ be a
vector lattice homomorphism, and let $x \in X$ with $0 \le x$.
Then $0 \le f x$.
\end{lemma}

\begin{lemma}
\label{lem:vecLatHom-map-posPart}
\lean{VecLatHom.map_posPart}
\leanok
\uses{def:vecLatHom}
Let $X$ and $Y$ be vector lattices and let $f \colon X \to Y$ be a
vector lattice homomorphism. Then $f x^{+} = (f x)^{+}$ for every
$x \in X$.
\end{lemma}

\begin{lemma}
\label{lem:vecLatHom-map-negPart}
\lean{VecLatHom.map_negPart}
\leanok
\uses{def:vecLatHom}
Let $X$ and $Y$ be vector lattices and let $f \colon X \to Y$ be a
vector lattice homomorphism. Then $f x^{-} = (f x)^{-}$ for every
$x \in X$.
\end{lemma}

\begin{lemma}
\label{lem:vecLatHom-map-disjoint}
\lean{VecLatHom.map_disjoint}
\leanok
\uses{def:vecLatHom, def:isVLDisjoint}
Let $X$ and $Y$ be vector lattices and let $f \colon X \to Y$ be a
vector lattice homomorphism. If $x, y \in X$ are disjoint, then
$f x$ and $f y$ are disjoint.
\end{lemma}

\begin{theorem}
\label{thm:vecLatHom-of-abs}
\lean{VecLatHom.of_abs, IsVecLatHom.of_abs}
\leanok
\uses{def:vecLatHom}
Let $X$ and $Y$ be vector lattices and let $f \colon X \to Y$ be a
real-linear map satisfying $f \lvert x \rvert = \lvert f x \rvert$
for every $x \in X$. Then $f$ is a vector lattice homomorphism.
\end{theorem}

\begin{theorem}
\label{thm:vecLatHom-of-disjoint}
\lean{IsVecLatHom.of_disjoint}
\leanok
\uses{def:vecLatHom, def:positive-operator, def:isVLDisjoint,
  lem:uniqueness-posPart}
Let $X$ and $Y$ be vector lattices and let $f \colon X \to Y$ be a
positive real-linear map. If $f$ preserves disjointness — that
is, $x \wedge y = 0$ implies $f x \wedge f y = 0$ for all $x, y
\in X$ — then $f$ is a vector lattice homomorphism.
\end{theorem}

\begin{lemma}
\label{lem:vecLatHom-symm}
\lean{VecLatHom.symm}
\leanok
\uses{def:vecLatHom}
Let $X$ and $Y$ be vector lattices and let $f \colon X \to Y$ be a
bijective vector lattice homomorphism. Then $f^{-1} \colon Y \to
X$ is also a vector lattice homomorphism.
\end{lemma}

\subsection{Vector lattice isomorphisms}

\begin{definition}
\label{def:vecLatEquiv}
\lean{VecLatEquiv}
\leanok
\uses{def:vecLatHom}
Let $X$ and $Y$ be vector lattices. A \emph{vector lattice
isomorphism} from $X$ to $Y$ is a real-linear bijection $e \colon
X \to Y$ that preserves the lattice operations $\vee$ and
$\wedge$.
\end{definition}

\begin{theorem}
\label{thm:positive-extensionEquiv}
\lean{Positive.extensionEquiv}
\leanok
\uses{def:vecLatEquiv, def:vector-lattice, def:is-vl-archimedean,
  thm:positive-extension, lem:positive-extension-nonneg,
  thm:positive-extension-positive, thm:positive-monotone-iff}
Let $X$ be a vector lattice, let $Y$ be an Archimedean vector
lattice, and let $\tau \colon X \to Y$ be a map satisfying $\tau x
\ge 0$ for every $x \ge 0$ in $X$ and $\tau (x + y) = \tau x +
\tau y$ for every $x, y \ge 0$ in $X$. If, in addition, $\tau$ is
injective and surjective on positive cones — that is, $\tau x =
\tau y$ implies $x = y$ for all $x, y \ge 0$ in $X$, and for every
$y \ge 0$ in $Y$ there exists $x \ge 0$ in $X$ with $\tau x = y$ —
then the linear extension $T x = \tau x^{+} - \tau x^{-}$ is a
vector lattice isomorphism $X \to Y$.
\end{theorem}

\begin{theorem}
\label{thm:vecLatEquiv-toContinuousLinearEquiv}
\lean{VecLatEquiv.toContinuousLinearEquiv}
\leanok
\uses{def:vecLatEquiv, def:banach-lattice,
  thm:positive-continuous, thm:positive-monotone-iff,
  lem:vecLatHom-monotone}
Let $X$ and $Y$ be Banach lattices. Every vector lattice
isomorphism $e \colon X \to Y$ is a continuous linear equivalence
$X \to Y$ — that is, both $e$ and $e^{-1}$ are continuous.
\end{theorem}

\subsection{Banach lattice isometries}

\begin{definition}
\label{def:banachLatEquiv}
\lean{BanachLatEquiv}
\leanok
\uses{def:vecLatEquiv, def:banach-lattice}
Let $X$ and $Y$ be Banach lattices. A \emph{Banach lattice
isometry} from $X$ to $Y$ is a real linear isometric equivalence
$e \colon X \to Y$ that preserves the lattice operations $\vee$
and $\wedge$.
\end{definition}

\section{The Riesz--Kantorovich theorem}

When the codomain $Y$ is a Dedekind complete vector lattice, the
space $\mathcal{L}_b(X, Y)$ of order bounded operators is itself a
Dedekind complete vector lattice. The lattice operations admit
explicit pointwise descriptions on the positive cone of $X$ — the
\emph{Riesz--Kantorovich formulas} — and there are companion
\emph{dual} formulas describing $f$ applied to a lattice expression
in $X$ as a supremum or infimum taken over operators dominated by a
fixed positive $f$. Finally, every order bounded operator from $X$
into a Dedekind complete vector lattice $Y$ is regular.

\subsection{The lattice structure on $\mathcal{L}_b(X, Y)$}

\begin{theorem}
\label{thm:orderBoundedHom-instLattice}
\lean{OrderBoundedHom.instLattice}
\leanok
\uses{def:orderBoundedHom, def:vector-lattice}
Let $X$ be a vector lattice and let $Y$ be a Dedekind complete
vector lattice. The space $\mathcal{L}_b(X, Y)$ of order bounded
operators is a lattice under the operator order.
\end{theorem}

\begin{theorem}
\label{thm:orderBoundedHom-instVectorLattice}
\lean{OrderBoundedHom.instVectorLattice}
\leanok
\uses{def:vector-lattice, def:orderBoundedHom,
  thm:orderBoundedHom-instLattice}
Let $X$ be a vector lattice and let $Y$ be a Dedekind complete
vector lattice. The space $\mathcal{L}_b(X, Y)$ of order bounded
operators is a vector lattice.
\end{theorem}

\subsection{Riesz--Kantorovich formulas}

\begin{theorem}
\label{thm:isLUB-posPart-apply}
\lean{OrderBoundedHom.isLUB_posPart_apply}
\leanok
\uses{def:orderBoundedHom, thm:orderBoundedHom-instVectorLattice,
  thm:riesz-decomposition}
Let $X$ be a vector lattice, let $Y$ be a Dedekind complete vector
lattice, let $f \in \mathcal{L}_b(X, Y)$, and let $x \in X$ with
$0 \le x$. Then $f^{+} x$ is the least upper bound of the set
$\{ f y : y \in X, \; 0 \le y \le x \}$.
\end{theorem}

\begin{theorem}
\label{thm:isLUB-negPart-apply}
\lean{OrderBoundedHom.isLUB_negPart_apply}
\leanok
\uses{def:orderBoundedHom, thm:orderBoundedHom-instVectorLattice,
  thm:isLUB-posPart-apply}
Let $X$ be a vector lattice, let $Y$ be a Dedekind complete vector
lattice, let $f \in \mathcal{L}_b(X, Y)$, and let $x \in X$ with
$0 \le x$. Then $f^{-} x$ is the least upper bound of the set
$\{ -f y : y \in X, \; 0 \le y \le x \}$.
\end{theorem}

\begin{theorem}
\label{thm:isLUB-sup-apply}
\lean{OrderBoundedHom.isLUB_sup_apply}
\leanok
\uses{def:orderBoundedHom, thm:orderBoundedHom-instVectorLattice,
  thm:isLUB-posPart-apply, lem:sup-eq-add-posPart}
Let $X$ be a vector lattice, let $Y$ be a Dedekind complete vector
lattice, let $f, g \in \mathcal{L}_b(X, Y)$, and let $x \in X$ with
$0 \le x$. Then $(f \vee g) x$ is the least upper bound of the set
$\{ f y + g z : y, z \in X, \; 0 \le y, \; 0 \le z, \; y + z = x \}$.
\end{theorem}

\begin{theorem}
\label{thm:isGLB-inf-apply}
\lean{OrderBoundedHom.isGLB_inf_apply}
\leanok
\uses{def:orderBoundedHom, thm:orderBoundedHom-instVectorLattice,
  thm:isLUB-sup-apply}
Let $X$ be a vector lattice, let $Y$ be a Dedekind complete vector
lattice, let $f, g \in \mathcal{L}_b(X, Y)$, and let $x \in X$ with
$0 \le x$. Then $(f \wedge g) x$ is the greatest lower bound of the
set $\{ f y + g z : y, z \in X, \; 0 \le y, \; 0 \le z, \; y + z =
x \}$.
\end{theorem}

\begin{theorem}
\label{thm:isLUB-abs-apply}
\lean{OrderBoundedHom.isLUB_abs_apply}
\leanok
\uses{def:orderBoundedHom, thm:orderBoundedHom-instVectorLattice,
  thm:isLUB-posPart-apply, thm:isLUB-negPart-apply,
  thm:isLUB-sup-apply}
Let $X$ be a vector lattice, let $Y$ be a Dedekind complete vector
lattice, let $f \in \mathcal{L}_b(X, Y)$, and let $x \in X$ with
$0 \le x$. Then $\lvert f \rvert x$ is the least upper bound of the
set $\{ \lvert f y \rvert : y \in X, \; \lvert y \rvert \le x \}$.
\end{theorem}

\subsection{Dual Riesz--Kantorovich formulas}

The dual formulas fix a positive operator $f \in \mathcal{L}_b(X,
Y)$ and describe $f$ applied to a lattice expression in $X$ as a
supremum or infimum taken over the order interval $[0, f]$ in
$\mathcal{L}_b(X, Y)$.

\begin{theorem}
\label{thm:isLUB-apply-posPart}
\lean{OrderBoundedHom.isLUB_apply_posPart}
\leanok
\uses{def:orderBoundedHom, thm:orderBoundedHom-instVectorLattice,
  thm:riesz-decomposition}
Let $X$ be a vector lattice, let $Y$ be a Dedekind complete vector
lattice, let $f \in \mathcal{L}_b(X, Y)$ with $0 \le f$, and let
$x \in X$. Then $f x^{+}$ is the least upper bound of the set
$\{ g x : g \in \mathcal{L}_b(X, Y), \; 0 \le g \le f \}$.
\end{theorem}

\begin{theorem}
\label{thm:isLUB-apply-negPart}
\lean{OrderBoundedHom.isLUB_apply_negPart}
\leanok
\uses{def:orderBoundedHom, thm:orderBoundedHom-instVectorLattice,
  thm:isLUB-apply-posPart}
Let $X$ be a vector lattice, let $Y$ be a Dedekind complete vector
lattice, let $f \in \mathcal{L}_b(X, Y)$ with $0 \le f$, and let
$x \in X$. Then $f x^{-}$ is the least upper bound of the set
$\{ -g x : g \in \mathcal{L}_b(X, Y), \; 0 \le g \le f \}$.
\end{theorem}

\begin{theorem}
\label{thm:isLUB-apply-sup}
\lean{OrderBoundedHom.isLUB_apply_sup}
\leanok
\uses{def:orderBoundedHom, thm:orderBoundedHom-instVectorLattice,
  thm:isLUB-apply-posPart}
Let $X$ be a vector lattice, let $Y$ be a Dedekind complete vector
lattice, let $f \in \mathcal{L}_b(X, Y)$ with $0 \le f$, and let
$x, y \in X$. Then $f (x \vee y)$ is the least upper bound of the
set $\{ g x + (f - g) y : g \in \mathcal{L}_b(X, Y), \; 0 \le g \le
f \}$.
\end{theorem}

\begin{theorem}
\label{thm:isGLB-apply-inf}
\lean{OrderBoundedHom.isGLB_apply_inf}
\leanok
\uses{def:orderBoundedHom, thm:orderBoundedHom-instVectorLattice,
  thm:isLUB-apply-sup}
Let $X$ be a vector lattice, let $Y$ be a Dedekind complete vector
lattice, let $f \in \mathcal{L}_b(X, Y)$ with $0 \le f$, and let
$x, y \in X$. Then $f (x \wedge y)$ is the greatest lower bound of
the set $\{ g x + (f - g) y : g \in \mathcal{L}_b(X, Y), \; 0 \le
g \le f \}$.
\end{theorem}

\begin{theorem}
\label{thm:isLUB-apply-abs}
\lean{OrderBoundedHom.isLUB_apply_abs}
\leanok
\uses{def:orderBoundedHom, thm:orderBoundedHom-instVectorLattice,
  thm:isLUB-apply-posPart}
Let $X$ be a vector lattice, let $Y$ be a Dedekind complete vector
lattice, let $f \in \mathcal{L}_b(X, Y)$ with $0 \le f$, and let
$x \in X$. Then $f \lvert x \rvert$ is the least upper bound of the
set $\{ g x - (f - g) x : g \in \mathcal{L}_b(X, Y), \; 0 \le g \le
f \}$.
\end{theorem}

\subsection{Dedekind completeness of $\mathcal{L}_b(X, Y)$}

\begin{theorem}
\label{thm:orderBoundedHom-isLUB-of-directedOn}
\lean{OrderBoundedHom.isLUB_of_directedOn}
\leanok
\uses{def:orderBoundedHom, thm:orderBoundedHom-instVectorLattice,
  thm:positive-extension, thm:positive-extension-positive,
  thm:positive-isOrderBounded}
Let $X$ be a vector lattice and let $Y$ be a Dedekind complete
vector lattice. Let $S \subseteq \mathcal{L}_b(X, Y)$ be a non-empty
upward-directed set which is bounded above. Then $S$ admits a least
upper bound $f \in \mathcal{L}_b(X, Y)$, and for every $x \in X$
with $0 \le x$ the value $f x$ is the least upper bound of the set
$\{ g x : g \in S \}$.
\end{theorem}

\begin{theorem}
\label{thm:orderBoundedHom-exists-isLUB}
\lean{OrderBoundedHom.exists_isLUB}
\leanok
\uses{def:orderBoundedHom, thm:orderBoundedHom-instVectorLattice,
  thm:orderBoundedHom-isLUB-of-directedOn}
Let $X$ be a vector lattice and let $Y$ be a Dedekind complete
vector lattice. Every non-empty bounded above subset of
$\mathcal{L}_b(X, Y)$ admits a least upper bound; that is,
$\mathcal{L}_b(X, Y)$ is Dedekind complete.
\end{theorem}

\subsection{Order bounded operators are regular}

\begin{theorem}
\label{thm:isOrderBounded-isRegularOp}
\lean{IsOrderBounded.isRegularOp}
\leanok
\uses{def:isOrderBounded, def:isRegularOp,
  thm:orderBoundedHom-instVectorLattice}
Let $X$ be a vector lattice and let $Y$ be a Dedekind complete
vector lattice. Every order bounded linear operator $f \colon X
\to Y$ is regular.
\end{theorem}

\section{The order dual and the strong dual}

This section specialises the operator theory of the previous
sections to scalar-valued functionals. The \emph{order dual} of a
vector lattice $X$ is the space $X^{\sim} = \mathcal{L}_b(X,
\mathbb{R})$ of order bounded real-valued linear functionals on
$X$. Since $\mathbb{R}$ is Dedekind complete, the
Riesz--Kantorovich machinery equips $X^{\sim}$ with a canonical
Dedekind complete vector lattice structure, and every element of
$X^{\sim}$ is regular. The \emph{strong dual} $X^{*}$ of a normed
vector lattice $X$ is the space of continuous real-valued linear
functionals. Every continuous functional is order bounded; on a
Banach lattice the converse also holds, so $X^{*}$ and $X^{\sim}$
coincide as sets, and the strong dual inherits the structure of a
Banach lattice with a solid lattice norm.

\subsection{The order dual}

\begin{definition}
\label{def:orderDualSpace}
\lean{OrderDualSpace}
\leanok
\uses{def:vector-lattice, def:orderBoundedHom}
Let $X$ be a vector lattice. The \emph{order dual} of $X$ is the
space $X^{\sim} = \mathcal{L}_b(X, \mathbb{R})$ of order bounded
real-valued linear functionals on $X$.
\end{definition}

\begin{lemma}
\label{lem:orderDualSpace-isRegularOp}
\lean{OrderDualSpace.isRegularOp}
\leanok
\uses{def:orderDualSpace, def:isRegularOp,
  thm:isOrderBounded-isRegularOp}
Let $X$ be a vector lattice. Every $\varphi \in X^{\sim}$ is
regular: there exist positive linear functionals $\varphi_1,
\varphi_2 \colon X \to \mathbb{R}$ such that $\varphi =
\varphi_1 - \varphi_2$.
\end{lemma}

\begin{definition}
\label{def:orderDualSpace-separatesPoints}
\lean{OrderDualSpace.SeparatesPoints}
\leanok
\uses{def:orderDualSpace}
The order dual $X^{\sim}$ \emph{separates points} of the vector
lattice $X$ if, for every $x \in X$, the condition $\varphi x = 0$
for all $\varphi \in X^{\sim}$ implies $x = 0$.
\end{definition}

\subsection{The strong dual of a normed vector lattice}

\begin{theorem}
\label{thm:strongDual-isOrderBounded}
\lean{StrongDual.isOrderBounded}
\leanok
\uses{def:isOrderBounded, def:normed-vector-lattice}
Let $X$ be a normed vector lattice. Every continuous real-linear
functional $\varphi \colon X \to \mathbb{R}$ is order bounded.
\end{theorem}

\begin{definition}
\label{def:strongDual-toOrderDualSpace}
\lean{StrongDual.toOrderDualSpace}
\leanok
\uses{def:orderDualSpace, def:normed-vector-lattice,
  thm:strongDual-isOrderBounded}
Let $X$ be a normed vector lattice. The map $X^{*} \to X^{\sim}$
sending a continuous functional $\varphi$ to itself, regarded as
an order bounded functional, is well defined.
\end{definition}

\subsection{The strong dual of a Banach lattice}

\begin{theorem}
\label{thm:strongDual-ofOrderDualSpace}
\lean{StrongDual.ofOrderDualSpace}
\leanok
\uses{def:orderDualSpace, def:banach-lattice,
  thm:isOrderBounded-continuous}
Let $X$ be a Banach lattice. Every order bounded real-linear
functional $\varphi \colon X \to \mathbb{R}$ is continuous.
\end{theorem}

\begin{theorem}
\label{thm:strongDual-equivOrderDualSpace}
\lean{StrongDual.equivOrderDualSpace}
\leanok
\uses{def:orderDualSpace, def:banach-lattice,
  thm:strongDual-isOrderBounded, thm:strongDual-ofOrderDualSpace}
Let $X$ be a Banach lattice. The strong dual $X^{*}$ and the order
dual $X^{\sim}$ coincide as real vector spaces: the inclusion
$X^{*} \to X^{\sim}$ is a real-linear bijection.
\end{theorem}

\begin{theorem}
\label{thm:strongDual-instLattice}
\lean{StrongDual.instLattice}
\leanok
\uses{def:banach-lattice, thm:strongDual-equivOrderDualSpace,
  thm:orderBoundedHom-instLattice}
Let $X$ be a Banach lattice. The strong dual $X^{*}$ is a lattice
under the partial order $\varphi \le \psi$ defined by $\varphi x
\le \psi x$ for every $x \in X$ with $0 \le x$, with lattice
operations transported from those of the order dual $X^{\sim}$
via the bijection $X^{*} \cong X^{\sim}$.
\end{theorem}

\begin{theorem}
\label{thm:strongDual-instVectorLattice}
\lean{StrongDual.instVectorLattice}
\leanok
\uses{def:vector-lattice, def:banach-lattice,
  thm:strongDual-instLattice,
  thm:orderBoundedHom-instVectorLattice}
Let $X$ be a Banach lattice. The strong dual $X^{*}$ is a vector
lattice.
\end{theorem}

\begin{theorem}
\label{thm:strongDual-instHasSolidNorm}
\lean{StrongDual.instHasSolidNorm}
\leanok
\uses{def:banach-lattice, thm:strongDual-instVectorLattice,
  thm:isLUB-abs-apply}
Let $X$ be a Banach lattice. The dual norm on $X^{*}$ is a solid
lattice norm: if $\lvert \varphi \rvert \le \lvert \psi \rvert$
in $X^{*}$, then $\lVert \varphi \rVert \le \lVert \psi \rVert$.
\end{theorem}

\begin{theorem}
\label{thm:strongDual-instNormedVectorLattice}
\lean{StrongDual.instNormedVectorLattice}
\leanok
\uses{def:normed-vector-lattice, def:banach-lattice,
  thm:strongDual-instVectorLattice,
  thm:strongDual-instHasSolidNorm}
Let $X$ be a Banach lattice. The strong dual $X^{*}$ is a normed
vector lattice under its dual norm.
\end{theorem}

\begin{theorem}
\label{thm:strongDual-instBanachLattice}
\lean{StrongDual.instBanachLattice}
\leanok
\uses{def:banach-lattice, thm:strongDual-instNormedVectorLattice}
Let $X$ be a Banach lattice. The strong dual $X^{*}$ is itself a
Banach lattice.
\end{theorem}

\begin{theorem}
\label{thm:strongDual-norm-abs}
\lean{StrongDual.norm_abs}
\leanok
\uses{def:banach-lattice, thm:strongDual-instHasSolidNorm}
Let $X$ be a Banach lattice. For every $\varphi \in X^{*}$,
$\lVert \, \lvert \varphi \rvert \, \rVert = \lVert \varphi
\rVert$.
\end{theorem}

\begin{theorem}
\label{thm:strongDual-norm-of-nonneg}
\lean{StrongDual.norm_of_nonneg}
\leanok
\uses{def:banach-lattice, thm:strongDual-instBanachLattice}
Let $X$ be a Banach lattice and let $\varphi \in X^{*}$ with
$0 \le \varphi$. Then
\[
  \lVert \varphi \rVert
  = \sup \{ \varphi x : x \in X, \; \lVert x \rVert \le 1, \;
    0 \le x \}.
\]
\end{theorem}

\begin{theorem}
\label{thm:strongDual-orderDual-separatesPoints}
\lean{StrongDual.orderDual_separatesPoints}
\leanok
\uses{def:orderDualSpace-separatesPoints, def:banach-lattice,
  thm:strongDual-equivOrderDualSpace}
Let $X$ be a Banach lattice. The order dual $X^{\sim}$ separates
points of $X$: for every $x \in X$, if $\varphi x = 0$ for all
$\varphi \in X^{\sim}$, then $x = 0$.
\end{theorem}

\section{The bidual of a Banach lattice}

The \emph{bidual} of a Banach lattice $X$ is the strong dual of
its strong dual, $X^{**} = (X^{*})^{*}$. Iterating the
construction of the previous section shows that $X^{**}$ is itself
a Banach lattice. The canonical evaluation map $X \to X^{**}$,
which embeds every normed space into its bidual as a linear
isometry, is in the lattice setting an isometric vector lattice
homomorphism: it preserves not only the norm but also the lattice
operations $\vee$, $\wedge$, and $\lvert \cdot \rvert$. In
particular, every Banach lattice sits inside its bidual as a
vector sublattice.

\begin{definition}
\label{def:bidualSpace}
\lean{BidualSpace}
\leanok
\uses{def:banach-lattice, thm:strongDual-instBanachLattice}
Let $X$ be a Banach lattice. The \emph{bidual} of $X$ is the
Banach lattice $X^{**} = (X^{*})^{*}$, the strong dual of the
strong dual of $X$.
\end{definition}

\begin{definition}
\label{def:bidualSpace-inclusion}
\lean{BidualSpace.inclusion}
\leanok
\uses{def:bidualSpace}
Let $X$ be a Banach lattice. The \emph{canonical embedding} of
$X$ into its bidual is the continuous real-linear map $X \to
X^{**}$ sending $x \in X$ to the evaluation functional $\varphi
\mapsto \varphi x$.
\end{definition}

\begin{theorem}
\label{thm:bidualSpace-norm-inclusion}
\lean{BidualSpace.norm_inclusion}
\leanok
\uses{def:bidualSpace-inclusion}
Let $X$ be a Banach lattice. The canonical embedding $X \to
X^{**}$ is a linear isometry: $\lVert \iota x \rVert = \lVert x
\rVert$ for every $x \in X$.
\end{theorem}

\begin{theorem}
\label{thm:bidualSpace-inclusion-nonneg}
\lean{BidualSpace.inclusion_nonneg}
\leanok
\uses{def:bidualSpace-inclusion, def:positive-operator}
Let $X$ be a Banach lattice. The canonical embedding $X \to
X^{**}$ is positive: if $0 \le x$ in $X$, then $0 \le \iota x$
in $X^{**}$.
\end{theorem}

\begin{theorem}
\label{thm:bidualSpace-inclusion-le-inclusion-iff}
\lean{BidualSpace.inclusion_le_inclusion_iff}
\leanok
\uses{def:bidualSpace-inclusion,
  thm:bidualSpace-inclusion-nonneg,
  thm:strongDual-orderDual-separatesPoints}
Let $X$ be a Banach lattice. For all $x, y \in X$, $\iota x \le
\iota y$ in $X^{**}$ if and only if $x \le y$.
\end{theorem}

\begin{theorem}
\label{thm:bidualSpace-inclusion-injective}
\lean{BidualSpace.inclusion_injective}
\leanok
\uses{def:bidualSpace-inclusion,
  thm:bidualSpace-norm-inclusion}
Let $X$ be a Banach lattice. The canonical embedding $X \to
X^{**}$ is injective.
\end{theorem}

\begin{theorem}
\label{thm:bidualSpace-inclusion-abs}
\lean{BidualSpace.inclusion_abs}
\leanok
\uses{def:bidualSpace-inclusion,
  thm:bidualSpace-inclusion-le-inclusion-iff,
  thm:isLUB-apply-posPart, thm:isLUB-posPart-apply}
Let $X$ be a Banach lattice. The canonical embedding $X \to
X^{**}$ preserves the modulus: $\iota \lvert x \rvert = \lvert
\iota x \rvert$ for every $x \in X$.
\end{theorem}

\begin{theorem}
\label{thm:bidualSpace-inclusion-sup}
\lean{BidualSpace.inclusion_sup}
\leanok
\uses{def:bidualSpace-inclusion,
  thm:bidualSpace-inclusion-abs, lem:sup-eq-add-posPart}
Let $X$ be a Banach lattice. The canonical embedding $X \to
X^{**}$ preserves binary suprema: $\iota (x \vee y) = \iota x
\vee \iota y$ for all $x, y \in X$.
\end{theorem}

\begin{theorem}
\label{thm:bidualSpace-inclusion-inf}
\lean{BidualSpace.inclusion_inf}
\leanok
\uses{def:bidualSpace-inclusion,
  thm:bidualSpace-inclusion-abs, lem:inf-eq-sub-posPart}
Let $X$ be a Banach lattice. The canonical embedding $X \to
X^{**}$ preserves binary infima: $\iota (x \wedge y) = \iota x
\wedge \iota y$ for all $x, y \in X$.
\end{theorem}

\begin{theorem}
\label{thm:bidualSpace-inclusionVecLatHom}
\lean{BidualSpace.inclusionVecLatHom}
\leanok
\uses{def:bidualSpace-inclusion, def:vecLatHom,
  thm:bidualSpace-inclusion-sup,
  thm:bidualSpace-inclusion-inf}
Let $X$ be a Banach lattice. The canonical embedding $X \to
X^{**}$ is a vector lattice homomorphism.
\end{theorem}
